\documentclass[10pt,letterpaper]{article}
\usepackage[top=0.85in,left=2.75in,footskip=0.75in]{geometry}

% amsmath and amssymb packages, useful for mathematical formulas and symbols
\usepackage{amsmath,amssymb}

% Use adjustwidth environment to exceed column width (see example table in text)
\usepackage{changepage}

% Use Unicode characters when possible
\usepackage[utf8x]{inputenc}

% textcomp package and marvosym package for additional characters
\usepackage{textcomp,marvosym}

% cite package, to clean up citations in the main text. Do not remove.
\usepackage{cite}

% Use nameref to cite supporting information files (see Supporting Information section for more info)
\usepackage{nameref,hyperref}

% line numbers
\usepackage[right]{lineno}

% ligatures disabled
\usepackage{microtype}
\DisableLigatures[f]{encoding = *, family = * }

% color can be used to apply background shading to table cells only
\usepackage[table]{xcolor}

% array package and thick rules for tables
\usepackage{array}

% enumerate package lets us use letters instead of numbers
\usepackage{enumerate}

% create "+" rule type for thick vertical lines
\newcolumntype{+}{!{\vrule width 2pt}}

% create \thickcline for thick horizontal lines of variable length
\newlength\savedwidth
\newcommand\thickcline[1]{%
  \noalign{\global\savedwidth\arrayrulewidth\global\arrayrulewidth 2pt}%
  \cline{#1}%
  \noalign{\vskip\arrayrulewidth}%
  \noalign{\global\arrayrulewidth\savedwidth}%
}

% \thickhline command for thick horizontal lines that span the table
\newcommand\thickhline{\noalign{\global\savedwidth\arrayrulewidth\global\arrayrulewidth 2pt}%
\hline
\noalign{\global\arrayrulewidth\savedwidth}}

\usepackage{color}

% Remove comment for double spacing
%\usepackage{setspace}
%\doublespacing

% Text layout
\raggedright
\setlength{\parindent}{0.5cm}
\textwidth 5.25in
\textheight 8.75in

% Bold the 'Figure #' in the caption and separate it from the title/caption with a period
% Captions will be left justified
\usepackage[aboveskip=1pt,labelfont=bf,labelsep=period,justification=raggedright,singlelinecheck=off]{caption}
\renewcommand{\figurename}{Fig}

% Use the PLoS provided BiBTeX style
\bibliographystyle{plos2015}

% Remove brackets from numbering in List of References
\makeatletter
\renewcommand{\@biblabel}[1]{\quad#1.}
\makeatother

% Leave date blank
\date{}

% Header and Footer with logo
\usepackage{lastpage,fancyhdr,graphicx}
\usepackage{epstopdf}
\pagestyle{myheadings}
\pagestyle{fancy}
\fancyhf{}
\setlength{\headheight}{27.023pt}
\lhead{\includegraphics[width=2.0in]{PLOS-submission.eps}}
\rfoot{\thepage/\pageref{LastPage}}
\renewcommand{\footrule}{\hrule height 2pt \vspace{2mm}}
\fancyheadoffset[L]{2.25in}
\fancyfootoffset[L]{2.25in}
\lfoot{\sf PLOS}


\newcommand{\rulemajor}[1]{\section*{#1}}

\begin{document}
\vspace*{0.2in}

\begin{flushleft}
{\Large
\textbf\newline{Ten Quick Tips for Making Things Findable}
}
\newline
\\
{Sarah Lin, MSLIS}\textsuperscript{1,*}
\\
\textbf{1} RStudio, Inc.
\\
\bigskip
* sarah.lin@rstudio.com
\end{flushleft}

\section*{Abstract}

Library science principles offer a framework and methodology for working through
the maze of information generated in professional life.  Information ecosystems
are composed of users, context, and content, and all three areas need to be
addressed to make information findable.  Scientists should define their users
and needs and then leverage the structural elements of the software they use to
create, store, and access information.  Only then can refine their content and
how they present it via naming and subject metadata.  Ultimately, improving
findability promotes scientists' work through contribution, iteration, and
dissemination of information and ideas.

\section*{Author summary}

Sarah Lin is the Information Architect and Digital Librarian at RStudio, PBC.

\section*{Introduction}

Researchers have always had to manage information of various kinds, but the
``exponential growth of electronic data increasingly calls for new means of
organizing and accessing information'' \cite{Hedden2016}. The resulting glut of
information, ``has also led to the creation of new technologies to help people
organize, find, and make better use of information''
\cite{Rosenfeld2015}. Trying to find a particular bit of information when needed
can waste precious time and cause frustration. On top of this, information often
resides in many formats and is shared with varied audiences, sometimes in
multiple versions, which makes it even harder to separate what's useful from
what's merely intriguing.

Luckily, librarians make a career out of helping people find information and use
information management technologies. Library science principles offer ways to
work through the maze of information generated in professional life, and
librarians' skills can be applied by any researcher who feels overwhelmed. The
ten quick tips in this paper are simple to understand, easy to apply, and
effective. They are based on the fact that information ecosystems have three
parts: users, context, and content \cite{Rosenfeld2015}. Solving the challenges
of information retrieval at the point of need therefore requires addressing the
people who use your information, the context within which it is created, and the
actual content of the information.

\rulemajor{1. Know Your Users}

The first step to making information findable is to establish who will be
finding information. In a professional context, your users are people who have a
``viable and legitimate interest in the work you're doing''
\cite{Covert2014}. Through contributions, they can help you iterate your own
work or expand upon it in their work. They are also transmitters of the
information you share through their own work and networks, and serving their
information needs ultimately benefits you as well as them.

Depending on your content and the nature of your role, you might have only a
small, expert audience. If your work is publicly available, however, accept the
idea that complete novices will find it, thereby making your actual user base
much larger than you might realize.

Novice users should be able to understand and navigate your content just as well
as subject experts. This much broader range of users means that you would expect
a wide spectrum of domain knowledge as well as facility navigating your
discipline's routine practices.

The information you wish to convey and the way it is currently organized
probably makes perfect sense to you. Its meaning for your users, however, is
determined by \emph{their} experience. That knowledge is ``\emph{whatever a user
  interprets} from the arrangement or sequence of things they encounter''
\cite{Covert2014}. This is another way of saying that the organizational
structure(s) you employ are a communication channel in their own right, in
addition to the domain knowledge you share. In addition to knowing who your
users are and how important they are to the impact of your work, remember that
any information you transmit to them is shaped by the mechanism(s) you use.

\rulemajor{2. Define Your Findability Problem}

Before you can design a solution, you must first establish both the extent of
the problem as well as the end state you seek. Those two elements together are
the \emph{why} that drives this process. Given how easy it is to create
information in a digital environment and the plethora of software and formats
you may employ personally and professionally, it is quite likely that your
information is a bit of a mess. ``The first step to taming any mess is to shine
a light on it so you can outline its edges and depths'' \cite{Covert2014}. This
means taking a look at both \emph{what} information you have and \emph{where} it
is located.  Once you know the full scope of the problem you have, you can move
towards greater findability.

Creating a picture of the end state you will achieve through information
organization can be personally motivating, yet in this context your users' needs
should be paramount. You must figure out what \emph{they} will consider a good
outcome, and ``[e]very decision you make should support what you've defined as
good'' \cite{Covert2014}. This may not completely overlap with what you
currently define as good for your own work, but the entire point of organizing
information is to make it findable to others as well as ourselves, so you must
balance your present workflow needs with the needs of your users, rather than
assuming that their needs perfectly match your own. And remember: your future
self is also one of your users, while your past self is effectively a novice, so
anything you do for others will likely pay off for yourself eventually.

Remember too that the question of how to find something can have several
different meanings. Users may need to find your work on the web, find a specific
item within a website, and/or find a particular piece of information within a
specific file or webpage. Depending on your content, you may have challenges in
all three areas. There are fine-tuning techniques appropriate for each; the
remaining tips related to context and content will help you address the full
scope of your findability problem.

Take, for example, the faculty website of a researcher coming up for tenure.

\begin{itemize}

\item
  The researcher created the website as a place to dump her vitae content and
  make it easier to fill out grant applications by having professional
  activities listed in one place.

\item
  A librarian might access that site looking for journal articles to include in
  the university's institutional repository, whereas a student might come
  looking for course information.

\item
  Tenure committee members might be interested in that laundry list of
  accomplishments, as a way to determine the impact factor of the researcher's
  published work. Three of these users need the researcher's list of
  professional publications, but three different `versions' of that information.

\item
  The researcher might only provide a plain text citation which meets her needs,
  but the librarian wouldn't get a link to the full text (or a PDF download) and
  the university administrator would not find each publication's impact factor.

\item
  The student needs are not considered at all.

\end{itemize}

Much of the content you create through your work has multiple audiences, so
figuring out the extent of \emph{who} needs to find \emph{what} is essential
before moving forward.

\rulemajor{3. Use Textual Structure}

Your information's context is one aspect of structural communication.
Findability issues at the document, post, or article level can be addressed by
taking advantage of the elements provided by each program you use to improve
organization and navigational support as well as retrieval support for your
users \cite{Hedden2016}.

For example, a key component of web searching is scanning text in search of
target information. Textual structure helps that process \cite{Krug2014}.
Examples of this are using formatted headers (rather than just enlarged text)
and bulleted or numbered lists, as well as \textbf{highlighting} terms that are
important. Depending on the software you are using, headings and table of
contents can be hyperlinked, which supports both scanning and navigation.

Textual content is created and aggregated in so many forms and using so many
different programs that specifying strategies beyond headings, lists, and
highlighting is difficult. While you want to exploit the features of the textual
program (your word processing software, data aggregator, or other
information-processing-tool) to the greatest degree possible, you must first
investigate each software for ways to do so and think about each program as a
way to notate information in addition to creating, manipulating, or storing
it. Examples include:

\begin{itemize}

\item
  GitHub offers tags that can be added to committ messages and are searchable.

\item
  Electronic lab notebooks of any kind can make use of an XML schema, like
  Darwin Core, EML, or FITS \cite{Briney2015}.

\item
  Using specific Google Docs heading levels creates a table of contents in real
  time, visible when the file is open.

\item
  Since CSV and spreadsheet files contain data, use a README file or create a
  data dictionary to provide information about the data and store it with raw
  data files.

\end{itemize}

\rulemajor{4. Add Metadata}

Just like people who end up with piles of photographs with nothing written on
the back, we all have digital mounds of files and content with no format or
subject information in the Properties (feature used by most modern operating
systems). Metadata, such as those contained in a file's Properties, adds extra
access points for information retrieval, which gives users more avenues to find
what they need. However, what metadata you are able to add depends on the
software used to create, store, and access your information, and on the file
formats that information is stored in.

Many file storage, word processing, and website construction programs have
built-in metadata capabilities, though they may be hard to find and harder to
understand how to leverage. The most difficult part may be getting into the
habit of adding metadata to your content after creation. You'll need to have
done some thinking about format and subject, but if your website software has
any metadata functionality (like tagging; see tip 9), you will want to take
advantage of the additional discoverability. If you ever have a choice of
software, one that makes use of metadata and structural labels enhances search
tremendously. It is also useful to examine how metadata can be transferred from
an old system to a new one if you have the luxury of multiple software options
(or have had a change forced on you).  For example:

\begin{itemize}

\item
  The tagging implemented on a WordPress website doesn't correspond to a
  taxonomy field in a \texttt{.yml} header file.

\item
  Keywords added to a journal article submission are not automatically added to
  the file metadata in PDF format, and pulling a citation from an article
  database doesn't give you the keywords as well.

\end{itemize}

Unfortunately, metadata is often software-specific and it is worth thinking
through what portability options you might have.  XML is an established standard
that would port metadata, if the software you're using accomodates it. Exporting
metadata in a portable file format, like plain text or XML, would have more
flexibility in migrating software formats.

\rulemajor{5. Use Search \emph{and} Browsing}

Research into information-seeking behavior shows that people use a combination
of searching \emph{and} browsing when they're trying to find information they
need. As they browse a website or document or file, they put together a mental
map of the content they could possibly find and search based on that map. ``In
the process, they modify their information requests as they learn more about
what they need and what information is available from the system''
\cite{Rosenfeld2015}. They refine their search further as they see results,
often navigating into a result or two to further their understanding. You have
probably seen or done something similar with a print book, trying to determine
if it's one you want by looking at the table of contents and the back
cover. These two functions work together because search allows users to find
information they know they need, whereas browsing allows users to find
information they don't know that they need \cite{Bates2002}.

Given this behavior, it's important to make information accessible both ways and
to ``look for ways to support moving easily from search to browse and back
again'' \cite{Rosenfeld2015}. Tags and other metadata help with searching, while
browsing a list of tags, a table of contents, or even a file directory tells
users a lot about the content contained in the information they are looking at.

Textual signifiers---elements within a text that convey meaning outside of the
words themselves---and software structures also help users navigate between
searching and browsing by making the information framework visible to the
users. That communication, ``enables the answers to users' questions to `rise'
to the surface'' and answer questions like, Where am I? What's here? Where can I
go from here?"  \cite{Rosenfeld2015}. These are essential wayfinding questions
on the path to findability and interact closely with the navigational choices
you make for your information.

A system of headings marked in a textual file is a great example of
intradocument wayfinding.  As a user scans the text, the headings act as topical
markers, so they can both summarize what information is contained in the
document, but also refine what they would search for based on the terms used in
those headings. Navigation bars on websites function in a similar way.  If the
user knows exactly what they are looking for, they can scan the menu and select
the option that matches their need.

In a situation where they don't know exactly what they need, the terms in the
menu help them understand both the vocabulary used in this domain as well as the
boundaries for what is included (i.e., terms are listed) or excluded (i.e., no
menu terms exist). In this way, search and browse work together to enhance
findability.  Though search might be thought of as a universal solution, without
knowing what it is possible to find and the words used to find it, search is not
a stand-alone solution.

\rulemajor{6. Mimic Real World Directions}

Even in a digital environment, our language mirrors that used for physical
directions. We ``visit'' or ``go to'' a website without actually changing our
physical location. Using the navigational metaphor completely is very helpful to
ensuring that you assist your users as much as possible, because ``the words you
use in the navigation systems and headings of [your content] help you find what
you're looking \emph{for}, but they also help you understand what you're looking
\emph{at}'' \cite{Arango2018}. Navigation is yet another structural element that
provides context to your users and communicates important information about your
work and how they can find what they need, so give your users a virtual map with
sign posts and directions so they can figure out where they need to go.

As a content creator, you might think that since your audience is educated, they
will have either an innate understanding of how you have structured your content
or unlimited patience to persist until they get what they need. In reality users
scan but don't read: they click on the first close thing they see and give up
very, very quickly. Menus, headings, file paths, breadcrumb trails on websites,
and tables of contents give users a sense of where the information resides as
part of the larger information ecosystem they are currently in as well as which
new paths they can take \cite{Krug2014}. Your directions should be consistent
throughout your information ecosystem just as highway signs are with regards to
appearance, style, and type of information,

\rulemajor{7. Name for Finding}

File and site naming affect findability both in terms of how they interact with
your search application and how human-readable they are once they are retrieved
as search results or through browsing. Naming can be a challenge because of the
need to balance a human-readable name that conveys the `aboutness' of the item
as well as a pattern or schema that is also machine-readable. For file
directories with many files, leveraging machine-readable components like sorting
with a particular date format can automatically place files in an order that
enhances findability. The utility of one choice over another hearkens back to
what you defined as good for your users.

There are many ways to develop a naming schema, largely related to the nature of
the information you create. At the most basic level, ``you should use consistent
names for the same reason that you use good file organization: so you can easily
find and use data later. Additionally, good naming helps you avoid duplicate
information'' \cite{Briney2015}. Researchers with multiple research projects or
significant complexity in their data sources might want to establish a unified
system of abbreviations for those projects or sources (these would be best
summarized in a data dictionary or README file). Consistency is key: choosing
your preferred case, preferred date format (YYYMMDD or YYYY-MM-DD will both sort
chronologically), and any project or source codes related to your research
and/or data types are all areas where standardization will benefit anyone
searching for information.

If you have things to name that are not files, such as projects, web pages, or
document headings, remember that the more generic the name the harder it is to
understand afterward, they can be nearly impossible to find things later if
their names do not reference their actual contents.  Try to ensure that your
names immediately and uniquely identify what they contain or reference. When
naming, you will want to imagine that someone will enter the name into a search
box at some point in the future. Will it lead them back to this particular item
or content? It should. Also, think about nicknames or shortened versions of your
names and make sure they are present in text or tags so that the content can be
discovered by a search engine and a user. Naming your raw data file ``raw'' or a
downloaded file ``download'' makes finding the information they contain nearly
impossible.  Renaming files can seem like a waste of precious time, but since
the research cycle doesn't end with publication \cite{Briney2015}, there is a
very high likelihood that someone will need to reuse your data and will have to
try to figure out what files corresponded to what part of your research.

\rulemajor{8. Use Tags}

Software metadata choices are part of the structural context of your
information, but what you choose to put in metadata fields related to content
subject and formats depends on your source material.  Additionally, digital
information can handle multiple tags, which aides searching.  This is a key
feature, because it means that you can now `file' information in multiple
locations--- something that was not possible in the era of print-only
information. Multiple tags also assist users from varied backgrounds because the
terms can be customized to each type of user your information has.

You need to be consistent in your depth of topical term assignment (how specific
your terms are) and selection of terms for subject and format (the number of
terms you use to describe each subject and format). Exactly which terms you
choose may not matter much for your discipline, beyond a need for consistency,
but in some disciplines the particular terms chosen matters a great
deal. However, if your discipline has an established thesaurus, it can be great
source for standardized subject terms. Thesauri have built-in subject
hierarchies that can help you create navigational structure. Examples include
the Astronomy Thesaurus, the Getty Art \& Architecture Thesaurus, the Education
Resources Information Center (ERIC) Thesaurus, the NASA Thesaurus, the UNESCO
Thesaurus, the GBA Thesaurus of Geosciences, and the Medical Subject Headings
(MeSH) \cite{ASI2020}.

These types of term lists are controlled vocabularies. They have a defined list
of ``official'' terms that is created and maintained by experts. There may or
may not be relationships built between terms, such as equivalencies (``CA'' for
``California''), broader/narrower terms (United States/California), and/or
replacement (weed USE marijuana). One benefit of using established subject terms
is that they that will match article databases and library catalogs that you and
your users might already be familiar with.

Alternatively, you can create your own taxonomy of subject keywords, which is
called a \emph{folksonomy}. Folksonomies are what you see with tags on Flickr
and Unsplash: the content creators assign terms as they see fit. Additionally,
users can tag with whatever terms they like and the list of possible terms is
both uncontrolled (no one is watching for synonyms, misspellings, or
inconsistencies) and crowd-sourced (created by non-experts). The good news is
that there is probably a place you can find tags in use that are a similar topic
or content to yours that will inspire the tag structure and/or terms you'll want
to use. If you create your own list of content terms, beware of plural/singular
discrepancies in your terms and be sure to review that list on a regular basis
for errors and misspellings.

\rulemajor{9. Understand the Difference Between Form and Substance}

When you create metadata tags--- whether controlled like a thesaurus or
uncontrolled like a folksonomy--- you need to address the distinction between
format and subject, which adds another layer of conceptual analysis to your
tagging process. Format describes what your content \emph{is}, while subject
describes what it is \emph{about} \cite{Joudrey2015}.  About-ness is the most
common content analysis, but you will likely have issues of is-ness that affect
your users and will want to add another layer of metadata to address it, where
appropriate. A simple example of this is a blog post on a website.  The post is
\emph{about} a subject, but it \emph{is} a blog post.  If it were your website,
you would want to tag that file both for its subject matter but also that it's a
blog post.  This enables users to search or filter by topic, format, or both.

Going back to your users, what subjects are important to them? And, do those
topics carry over or change between differences in format? This is basically a
question of combined terms: are your format terms uniquely matched to topical
terms (blog posts are always about news) or do you have multiple topics in each
format (blog posts and tutorials on the same subject)? When you think of file
formats, sometimes the software chooses the format which in turn signals what
the content is, for example, data stored in CSV files or a presentation saved as
a Keynote or PowerPoint file.

In these cases, you know a CSV file will contain some type of a data to be
analyzed, whereas Keynote and PowerPoint will have slides that are viewed
unidirectionally. However, dissemination sometimes changes a file's format
(e.g., printing slides to a PDF), in which case careful naming and metadata
conventions would recommend putting the format into the name and/or file
metadata. Once you are aware of format and subject analysis, you will likely see
more opportunities to distinguish information sources for your users. You might
create a folder specifically for conference presentations that could include
both .ppt and .pdf file extensions, where each presentation's files have the
exact same name (with different filetypes).  Likewise, journal articles you
store will need a naming convention that distinguishes between articles you have
written and those relied on for research.

\rulemajor{10. Do Not Abbrvt}

As an expert in your field, you ``can no longer imagine what it's like to
\emph{not} see the world that way'' \cite{Wilson2020}, and have to work to
ensure that the information you share with your users is accessible to
them. Spelling out acronyms and abbreviations (or hyperlinking to definitions)
breaks down the exclusivity of language used by a select group and makes
newcomers feel welcome.

Remember that acronyms are often repurposed by different professions or
disciplines: what seems obvious to you is probably not obvious to a number of
other people. Given that your discipline likely has many common abbreviations,
you can write out all abbreviations and acronyms in their full form, write them
out the first time they appear, or create or point to a term dictionary.

\subsection{Conclusion}

Changing work habits is hard, and you must decide whether to start new work with
a clean slate of new information principles or if you must go back and
alter---perhaps rename---existing work products \cite{Briney2015}.  Remember
that while ``perfection isn't possible, progress is'' \cite{Covert2014}.
Creating and structuring good information for your users is your goal, but these
tips and the ``ways you enforce your way of doing things changes how users think
about the place[s] you made and perhaps ultimately, how they think about you''
\cite{Covert2014}. Circling back to the network of users you are situated among,
making your stuff findable is a truly impactful way to disseminate your work to
them that will positively change your network and your career.

\bibliography{10-findable}

\end{document}
